\documentclass[11pt]{amsart}
\usepackage[T1]{fontenc}
\usepackage{amsmath,amssymb,amsthm}
\usepackage[colorlinks=true,linkcolor=blue,citecolor=blue,urlcolor=blue]{hyperref}

\newtheorem*{theorem*}{Partial theorem}

\newcommand{\Mass}{\mathbb M}
\newcommand{\R}{\mathbb R}
\newcommand{\Lip}{\operatorname{Lip}}

\begin{document}

\title[Finite-Dimensional Banach Transfer]{Finite-Dimensional Banach Transfer for the Alberti--Marchese Flat-Chain Theorem}
\author{Run fa\_banach\_001}
\date{2026-06-19}
\maketitle

\noindent
\textbf{Status.} \texttt{partial\_result\_likely\_valid}. This packet records a
finite-dimensional Banach-space subcase adjacent to the higher-dimensional
Banach extension question noted in Bate--Caputo--Takac--Valentine--Wald,
arXiv:2508.08017.

\section*{Source Target}

The source paper proves the \(1\)-dimensional flat-chain conjecture and a
Banach-space version of the Alberti--Marchese structure theorem for
\(1\)-currents. In its discussion of higher-dimensional currents, the authors
state that for \(k\ge 2\) the strict polyhedral approximation result fails in
every infinite-dimensional Hilbert space and that it is therefore unclear
whether Alberti--Marchese's Theorem 1.1 can be extended beyond
finite-dimensional Euclidean space.

The supporting Euclidean theorem is Alberti--Marchese, \emph{On the structure
of flat chains with finite mass}, arXiv:2311.06099, Theorem 1.1: every
real \(k\)-dimensional flat chain with finite mass in \(\R^d\), \(1\le k<d\),
is the restriction of a boundaryless normal \(k\)-current to a Borel set, with
mass at most \((2+\varepsilon)\) times the mass of the original chain.

\section*{Proof Intuition}

Finite-dimensional Banach spaces are linearly bi-Lipschitz to Euclidean space.
The structure theorem can therefore be transported across such an isomorphism:
push the flat chain to \(\R^d\), use Alberti--Marchese, and pull the resulting
normal current back. The price is exactly the usual Lipschitz mass distortion:
a \(k\)-current mass changes by at most the \(k\)-th power of the Lipschitz
constant.

\begin{theorem*}
Let \(E\) be a \(d\)-dimensional Banach space, let \(1\le k<d\), and let
\(T\) be a real \(k\)-dimensional flat chain with finite mass in \(E\).
Let \(A:E\to \R^d\) be a linear isomorphism, where \(\R^d\) carries its
Euclidean norm. Then for every \(\varepsilon>0\) there exist a boundaryless
normal \(k\)-current \(N\) in \(E\) and a Borel set \(B\subset E\) such that
\[
  T=N\llcorner B
\]
and
\[
  \Mass_E(N)
  \le (2+\varepsilon)\,\|A\|^k\|A^{-1}\|^k\,\Mass_E(T).
\]
Consequently the constant may be written with the \(k\)-th power of the
Banach--Mazur distance from \(E\) to \(\ell_2^d\), up to arbitrarily small
slack in the choice of \(A\).
\end{theorem*}

\begin{proof}
Because \(A\) is bi-Lipschitz, pushforward by \(A\) sends flat \(k\)-chains in
\(E\) to flat \(k\)-chains in \(\R^d\). Indeed, if \(P_i\) is a polyhedral
flat-norm approximating sequence for \(T\), then \(A_\#P_i\) is a polyhedral
approximating sequence for \(A_\#T\), with flat distances controlled by the
Lipschitz constants of \(A\). Also
\[
  \Mass_{\R^d}(A_\#T)\le \|A\|^k \Mass_E(T).
\]

Apply Alberti--Marchese's Euclidean theorem to \(A_\#T\). For any
\(\delta>0\) there are a boundaryless normal \(k\)-current \(N'\) in
\(\R^d\) and a Borel set \(B'\subset\R^d\) such that
\[
  A_\#T=N'\llcorner B'
\]
and
\[
  \Mass_{\R^d}(N')\le (2+\delta)\Mass_{\R^d}(A_\#T).
\]

Set \(N:=(A^{-1})_\#N'\) and \(B:=A^{-1}(B')\). Pushforward commutes with
boundary, so \(\partial N=(A^{-1})_\#\partial N'=0\), and \(N\) is normal
because normality is preserved under Lipschitz pushforward. Restriction to a
Borel set commutes with pushforward by the bijection \(A^{-1}\), hence
\[
  N\llcorner B
  =(A^{-1})_\#(N'\llcorner B')
  =(A^{-1})_\#A_\#T
  =T.
\]
Finally,
\[
  \Mass_E(N)
  \le \|A^{-1}\|^k \Mass_{\R^d}(N')
  \le (2+\delta)\|A^{-1}\|^k\|A\|^k\Mass_E(T).
\]
Taking \(\delta=\varepsilon\) gives the displayed estimate. Taking the
infimum over linear isomorphisms \(A:E\to\ell_2^d\), with the usual small
slack because the infimum need not be attained, gives the Banach--Mazur
formulation.
\end{proof}

\section*{Verification and Limitations}

The argument uses only standard functorial properties of metric currents and
flat chains under Lipschitz maps, plus Alberti--Marchese's Euclidean theorem.
No computation is involved.

This is not a solution of the infinite-dimensional Banach problem mentioned in
the source paper. It also does not prove the sharp \((2+\varepsilon)\) mass
constant for non-Euclidean finite-dimensional norms, and it does not identify
arbitrary metric \(k\)-currents in Banach spaces with flat chains.

\section*{Search Notes}

The run indexes were searched for arXiv:2508.08017, the source title,
Alberti--Marchese, flat-chain conjecture, finite-dimensional Banach spaces,
metric currents, and normal-current restriction language. Web searches on
June 19, 2026 for the same finite-dimensional Banach transfer phrases found
the Alberti--Marchese paper and adjacent metric-current literature, but no
direct finite-dimensional Banach statement.

\section*{References}

\begin{itemize}
\item D. Bate, E. Caputo, J. Takac, P. Valentine, P. Wald,
\emph{Structure of Metric \(1\)-currents: approximation by normal currents and
representation results}, arXiv:2508.08017.
\item G. Alberti and A. Marchese, \emph{On the structure of flat chains with
finite mass}, arXiv:2311.06099.
\end{itemize}

\end{document}
